% =============================================================================
%  CHAPTER 14: Custom Commands & Environments
% =============================================================================

\chapter{Custom Commands \& Environments}

\section{New Commands (\texttt{\textbackslash newcommand})}

\begin{lstlisting}[language={[LaTeX]TeX}, caption=newcommand syntax]
% Basic: no arguments
\newcommand{\myname}{John Doe}
Hello, I am \myname.          % Output: Hello, I am John Doe.

% With arguments (1 required argument):
\newcommand{\greet}[1]{Hello, #1!}
\greet{Alice}                 % Output: Hello, Alice!

% With multiple arguments:
\newcommand{\add}[2]{#1 + #2}
$\add{a}{b}$                  % Output: a + b

% With optional arguments:
\newcommand{\hello}[2][World]{Hello, #2 from #1!}
\hello{Alice}                 % Output: Hello, Alice from World!
\hello[Bob]{Alice}            % Output: Hello, Alice from Bob!

% Maximum 9 arguments (#1 through #9)
\end{lstlisting}

\subsection{Practical Examples}

\begin{lstlisting}[language={[LaTeX]TeX}, caption=Practical custom commands]
% Shorthand for common formatting:
\newcommand{\pkg}[1]{\texttt{#1}}
\newcommand{\cmd}[1]{\texttt{\textbackslash#1}}

% Math shortcuts:
\newcommand{\R}{\mathbb{R}}              % Real numbers
\newcommand{\norm}[1]{\left\|#1\right\|} % Norm notation
\newcommand{\abs}[1]{\left|#1\right|}    % Absolute value
\newcommand{\inner}[2]{\langle #1, #2 \rangle} % Inner product
\newcommand{\dd}{\,\mathrm{d}}           % Differential d

% Usage:
% $f\colon \R^n \to \R$
% $\norm{x} = \sqrt{\inner{x}{x}}$
% $\int_0^1 f(x) \dd x$
\end{lstlisting}

Results:
\[ f\colon \mathbb{R}^n \to \mathbb{R}, \quad \|x\| = \sqrt{\langle x, x \rangle} \]

% ------------------------------------------------------------------
\section{Renaming Commands (\texttt{\textbackslash renewcommand})}

\begin{lstlisting}[language={[LaTeX]TeX}, caption=renewcommand]
% Override existing commands:
\renewcommand{\contentsname}{Contents}        % Rename TOC title
\renewcommand{\figurename}{Fig.}              % Rename figure label
\renewcommand{\tablename}{Table}              % Rename table label
\renewcommand{\bibname}{References}           % Rename bibliography

% Check if command exists:
\providecommand{\mycommand}{default value}   % Define only if not already
\end{lstlisting}

% ------------------------------------------------------------------
\section{New Environments}

You can define entire environments (begin/end pairs) with \verb|\newenvironment|:

The syntax is:
\verb|\newenvironment{name}[args]{begin-code}{end-code}|

\textbf{Example 1} -- a simple warning box:
\begin{lstlisting}[language={[LaTeX]TeX}]
\newenvironment{warning}
    {\begin{center}\bfseries Warning: }
    {\end{center}}
\end{lstlisting}

\textbf{Example 2} -- a box with a title argument:
\begin{lstlisting}[language={[LaTeX]TeX}]
\newenvironment{infobox}[1]
    {\begin{center}\fbox{\parbox{0.9\textwidth}%
     {\textbf{##1}\par\smallskip}}}
    {\end{center}}
\end{lstlisting}

% ------------------------------------------------------------------
\section{Lengths and Counters}

\subsection{Defining Lengths}

\begin{Verbatim}[frame=single, fontsize=\small, label=Custom lengths]
% Create a new length:
\newlength{\mywidth}
\setlength{\mywidth}{5cm}

% Use it:
\hspace{\mywidth}

% Arithmetic on lengths:
\addtolength{\mywidth}{2cm}           % Add 2cm
\settowidth{\mywidth}{Hello World}    % Set to width of text
\settoheight{\myheight}{Text}         % Set to height
\settodepth{\mydepth}{Text}           % Set to depth
\end{Verbatim}

\subsection{Standard Lengths}

\begin{table}[H]
\centering
\begin{tabular}{ll}
\toprule
\textbf{Length} & \textbf{Controls} \\
\midrule
\verb|\textwidth|      & Width of text area \\
\verb|\linewidth|      & Width of current line \\
\verb|\columnwidth|    & Width of column \\
\verb|\paperwidth|     & Width of paper \\
\verb|\paperheight|    & Height of paper \\
\verb|\parindent|      & Paragraph indentation \\
\verb|\parskip|        & Space between paragraphs \\
\verb|\baselineskip|   & Distance between baselines \\
\verb|\columnsep|      & Space between columns \\
\bottomrule
\end{tabular}
\caption{Standard \latex{} lengths}
\end{table}

\subsection{Counters}

\begin{lstlisting}[language={[LaTeX]TeX}, caption=Counters]
% Define a new counter:
\newcounter{mycounter}

% Set, step, and use:
\setcounter{mycounter}{0}
\stepcounter{mycounter}      % Increment by 1
\addtocounter{mycounter}{5}  % Add 5

% Display:
\arabic{mycounter}   % 1, 2, 3
\roman{mycounter}    % i, ii, iii
\Roman{mycounter}    % I, II, III
\alph{mycounter}     % a, b, c
\Alph{mycounter}     % A, B, C
\fnsymbol{mycounter} % *, dagger, double-dagger, etc.
\end{lstlisting}

% ------------------------------------------------------------------
\section{Conditional Commands}

\begin{lstlisting}[language={[LaTeX]TeX}, caption=Conditionals]
% ifthen package:
\usepackage{ifthen}

\newcommand{\checkodd}[1]{%
    \ifthenelse{\isodd{#1}}%
        {#1 is odd}%
        {#1 is even}%
}

% Boolean flags:
\newboolean{draftmode}
\setboolean{draftmode}{true}

\ifthenelse{\boolean{draftmode}}%
    {\watermark{DRAFT}}%
    {}
\end{lstlisting}

% ------------------------------------------------------------------
\section{Boxes}

\begin{Verbatim}[frame=single, fontsize=\small, label=Boxes]
% mbox: prevents line breaks
\mbox{This text stays on one line}

% fbox: framed box
\fbox{Framed text}

% makebox: custom box with alignment
\makebox[5cm][c]{Centered in 5cm}
% Alignment: l=left, c=center, r=right, s=stretch

% parbox: paragraph in a box
\parbox{3cm}{This text wraps inside a 3cm-wide box.}

% minipage: like parbox but more powerful
\begin{minipage}[t]{0.45\textwidth}
    Left column content.
\end{minipage}
\hfill
\begin{minipage}[t]{0.45\textwidth}
    Right column content.
\end{minipage}

% raisebox: vertically offset text
\raisebox{5pt}{Raised text}
\raisebox{-5pt}{Lowered text}
\end{Verbatim}

Side-by-side minipage example:

\begin{minipage}[t]{0.45\textwidth}
    \textbf{Left Column}\\
    This content sits in the left minipage. It wraps automatically within the specified width.
\end{minipage}%
\hfill
\begin{minipage}[t]{0.45\textwidth}
    \textbf{Right Column}\\
    This content sits in the right minipage. Using minipages we can create multi-column layouts anywhere.
\end{minipage}
